\begin{figure}[htb]
     \centering
    \scalebox{1.0}{
% diamand in der mitte
% \begin{tikzpicture}
%  \tikzset{every state/.style={minimum size=1pt},>=stealth'}
%  \tikzstyle{dotted}=                  [dash pattern=on \pgflinewidth off 2pt]
   
 \begin{tikzpicture}[
 	>=latex',
 	shorten >=1pt,
 	node distance=6em,
 	on grid,
 	auto
 	]
 	\tikzset{every state/.style={minimum size=1pt},
 		mybox/.style={
 			draw=red,
 			dashed,
 			rounded corners,
 			inner sep=13pt
 		}
 	}
 	
 	%---------------- Left automaton ----------------%
 	
 	\node[state] (2) {};
 	\node[state] (3) [below left=of 2] {};
 	\node[state] (4) [below right=of 2] {};
 	\node[state] (5) [left=of 3] {};
 	\node[state] (6) [below right=of 3] {};
 	\node[state] (8) [above right=of 4] {};
 	\node[state] (9) [above=of 8] {};
 	
 	\path[->]
 	(2) edge node {$a$} (3)
 	(3) edge[bend left] node {$a$} (5)
 	(5) edge[bend left] node {$a$} (3)
 	(3) edge node {$b$} (6)
 	(6) edge[loop below] node {$a,b$} (6)
 	(5) edge[bend right=70] node {$b$} (6);
 	
 	\path[->]
 	(2) edge node {$b$} (4)
 	(4) edge[loop right] node {$b$} (4)
 	(4) edge node {$a$} (6);
 	
 	\path[->]
 	(8) edge[bend left] node {$b$} (9)
 	(9) edge[bend left] node {$b$} (8)
 	(8) edge node {$a$} (4)
 	(9) edge[bend right] node {$a$} (4);
 	
 	%---------------- Arrow ----------------%
 	
 	\node (P) [right=6em of 4] {\LARGE $\Rightarrow$};
 	
 	%---------------- Right automaton ----------------%
 	
 	\node[state] (S1) [right=3em of P] {};
 	\node[state] (S2) [above right=of S1] {};
 	\node[state] (S3) [below right=of S2] {};
 	\node[state] (S4) [below left=of S3] {};
 	\node[state] (S5) [above right=of S3] {};
 	
 	\path[->]
 	(S2) edge node {$a$} (S1)
 	(S2) edge node {$b$} (S3)
 	(S3) edge node {$a$} (S4)
 	(S1) edge node {$b$} (S4)
 	(S5) edge node {$a$} (S3);
 	
 	\path[->]
 	(S1) edge[loop above] node {$a$} (S1)
 	(S3) edge[loop right] node {$b$} (S3)
 	(S5) edge[loop above] node {$a,b$} (S5)
 	(S4) edge[loop below] node {$a,b$} (S4);
 	
 	%---------------- Background rectangles ----------------%
 	
 	\begin{pgfonlayer}{background}
 		
 		% States 3 and 5
 		\node[mybox, fit=(3)(5)] {};
 		
 		% State 6
 		\node[
 		mybox,
 		fit=(6), inner sep=8
 		] {};
 		
 		% State 2
 		\node[
 		mybox,
 		fit=(2), inner sep=8
 		] {};
 		
 		% State 4
 		\node[
 		mybox,
 		fit=(4), inner sep=8
 		] {};
 		
 		% States 8 and 9
 		\node[
 		mybox,
 		fit=(8)(9)
 		] {};
 		
 	\end{pgfonlayer}
 	
 \end{tikzpicture}
}
  \caption[Example quotient automaton of \texorpdfstring{$\rho$}{rho}-automaton]{
  	Left: A $\rho$-automaton $\mathcal A$ (that is commutative and synchronizing). The strongly connected components of $\mathcal A$
  	are surrounded by boxes. Right: The automaton $\mathcal A / {\eqreach}$.}
  %    have been framed by red boxes.}
  \label{fig:ex_comm_aut}
\end{figure}





%With Theorem~\ref{sync_max:thm:com_waa_same_sync_words}, 
%a commutative automaton with $n$
%states can only be sync-maximal when $2^n - n \le n$.
%This leaves only the possibilities $n \in \{1,2\}$.
%As sync-maximality implies that it is synchronizing \todo{das irgendwo hinschreiben als claim oder so?}
%the claim follows with Proposition~\ref{prop:two_state_comm_sync}.
%Conversely, that the two types of automata are commutative
%and sync-maximal is easily seen.



%\begin{toappendix2} 
%
%\begin{figure}[htb]
%     \centering
%    \scalebox{.7}{
%% diamand in der mitte
%% \begin{tikzpicture}
%%  \tikzset{every state/.style={minimum size=1pt},>=stealth'}
%%  \tikzstyle{dotted}=                  [dash pattern=on \pgflinewidth off 2pt]
%   
%  \begin{tikzpicture}[>=latex',shorten >=1pt,node distance=2cm,on grid,auto]
% \tikzset{every state/.style={minimum size=1pt}}
%    \node[state] (10) at (11,2) {};
%   \node[state] (11) at (14,1) {}; %[below right of=10] {};
%   \node[state] (12) at (10.5,-0.5) {}; % [below of=11] {};
%   \node[state] (13) [below of=12] {};
%   \node[state] (14) [right of=13] {};
%   \node[state] (15) [left of=11] {};
%   \node[state] (16) [right of=12] {};
%   
%   \path[->] (10) edge [bend left] node {$a$} (11)
%             (11) edge node {$b$} (16)
%             (10) edge [left] node {$b$} (12)
%            % (11) edge [loop left] node {$a$} (11)
%             (12) edge [left] node {$b$} (13)
%             (13) edge [loop below] node {$b$} (12)
%             (14) edge [loop below] node {$b$} (14)
%             (13) edge [bend left] node {$a$} (14)
%             (14) edge [bend left] node {$a$} (13);
%    
% % \path[->] (12) edge [loop right] node {$a$} (12);
% 
%   \path[->] (11) edge [bend left] node {$a$} (15)
%             (15) edge [bend left] node {$a$} (11);
%             
%   \path[->] (16) edge [bend left] node {$a$} (12)
%             (12) edge [bend left] node {$a$} (16);
%             
%   \path[->] (15) edge [left,pos=.3] node {$b$} (12);
%   \path[->] (16) edge node {$b$} (14);
%   
%\end{tikzpicture}
%}
%   
%   \caption{A commutative automaton that is not synchronizing.}
%  \label{fig:ex_comm_aut2}
%\end{figure}
%\end{toappendix2}